\section{Introduction}\label{sec:introduction}

Riemann's rearrangement theorem states that every conditionally
convergent real series has rearrangements converging to any prescribed
element of $\mathbb R\cup\{-\infty,+\infty\}$, as well as
rearrangements that diverge by oscillation \cite{Riemann}.  This
flexibility leads naturally to the following question: how many
permutations are required to disrupt every conditionally convergent
series?

To the best of the author's knowledge, this question first appeared
explicitly in a 2015 MathOverflow post by Hardy \cite{HardyMO}.  Blass,
Brendle, Brian, Hamkins, Hardy, and Larson subsequently introduced the
corresponding cardinal characteristic $\rr$, called the
\emph{rearrangement number}, and systematically studied it
\cite{BBBHHL}.  They proved
\[
 \max\{\covN,\bb\}\leq\rr\leq\nonM.
\]
The definitions of these cardinals are recalled in
Section~\ref{sec:preliminaries}.

The relevant characteristics occupy adjacent levels of Cichoń's
diagram, as illustrated in Figure~\ref{fig:intro-cichon}.  The two
lower bounds are incomparable in ZFC: both $\covN<\bb$ and
$\bb<\covN$ are consistent.

\begin{figure}[ht]
 \centering
 \begin{tikzpicture}[
   x=1cm,
   y=1cm,
   cardinal/.style={font=\small,inner sep=2pt},
   classical/.style={->,thin},
   rrbound/.style={->,thick,densely dashed},
   rrnode/.style={font=\large,draw,rounded corners,inner sep=4pt}
  ]
  \node[cardinal] (covN) at (0,2.8) {$\operatorname{cov}(\cN)$};
  \node[cardinal] (nonM) at (3,2.8) {$\operatorname{non}(\cM)$};
  \node[cardinal] (cofM) at (6,2.8) {$\operatorname{cof}(\cM)$};
  \node[cardinal] (cofN) at (9,2.8) {$\operatorname{cof}(\cN)$};
  \node[cardinal] (c) at (11.5,2.8) {$\mathfrak c$};

  \node[cardinal] (b) at (3,1.4) {$\bb$};
  \node[cardinal] (d) at (6,1.4) {$\dd$};

  \node[cardinal] (aleph1) at (-2.5,0) {$\aleph_1$};
  \node[cardinal] (addN) at (0,0) {$\operatorname{add}(\cN)$};
  \node[cardinal] (addM) at (3,0) {$\operatorname{add}(\cM)$};
  \node[cardinal] (covM) at (6,0) {$\operatorname{cov}(\cM)$};
  \node[cardinal] (nonN) at (9,0) {$\operatorname{non}(\cN)$};

  \draw[classical] (covN) -- (nonM);
  \draw[classical] (nonM) -- (cofM);
  \draw[classical] (cofM) -- (cofN);
  \draw[classical] (addN) -- (addM);
  \draw[classical] (addM) -- (covM);
  \draw[classical] (covM) -- (nonN);
  \draw[classical] (addN) -- (covN);
  \draw[classical] (addM) -- (b);
  \draw[classical] (b) -- (nonM);
  \draw[classical] (covM) -- (d);
  \draw[classical] (d) -- (cofM);
  \draw[classical] (nonN) -- (cofN);
  \draw[classical] (b) -- (d);
  \draw[classical] (aleph1) -- (addN);
  \draw[classical] (cofN) -- (c);

  \node[rrnode] (rr) at (1.55,1.95) {$\rr$};
  \draw[rrbound] (covN) -- (rr);
  \draw[rrbound] (b) -- (rr);
  \draw[rrbound] (rr) -- (nonM);
 \end{tikzpicture}
 \caption{Cichoń's diagram with the previously known bounds for $\rr$.}
 \label{fig:intro-cichon}
\end{figure}

Explicitly, they asked the following.

\begin{question}[\cite{BBBHHL}, Question~46]\label{q:intro-rr-nonM}
Does ZFC prove that $\rr=\nonM$?
\end{question}

The question was repeated as Problem~10 by Brech, Brendle, and Telles
\cite{BrechBrendleTelles}.
In this paper, we settle this question by proving that the answer is positive.

\begin{theorem}[Main Theorem]\label{thm:main}
In ZFC, $\rr=\nonM$.
\end{theorem}
